\documentclass[11pt]{article}
\usepackage[margin=1in]{geometry}
\usepackage[T1]{fontenc}
\usepackage[utf8]{inputenc}
\usepackage{lmodern}
\usepackage{microtype}
\usepackage{array}
\usepackage{booktabs}
\usepackage{enumitem}
\usepackage{titlesec}
\usepackage{xcolor}
\usepackage{hyperref}
\definecolor{bpsink}{HTML}{1C2530}
\definecolor{bpsaccent}{HTML}{2534B8}
\definecolor{bpsmuted}{HTML}{5B6675}
\definecolor{bpsenh}{HTML}{7A4BD0}
\definecolor{bpsline}{HTML}{CDD3DB}
\definecolor{bpssoft}{HTML}{F0F2F5}
\definecolor{bpswarnbg}{HTML}{FDF6E6}
\definecolor{bpswarnink}{HTML}{7A5A00}
\hypersetup{
  colorlinks=true,
  linkcolor=bpsaccent,
  urlcolor=bpsaccent,
  pdftitle={BPS Coding Scheme Dossier}
}
\setlist[itemize]{leftmargin=1.5em,itemsep=2pt,topsep=3pt}
\setlist[description]{style=nextline,leftmargin=0em,labelsep=0.6em,itemsep=4pt}
\titleformat{\section}{\color{bpsink}\large\bfseries}{}{0em}{}[{\color{bpsline}\titlerule}]
\titlespacing*{\section}{0pt}{1.1em}{0.5em}
\newcommand{\field}[1]{\nolinkurl{#1}}
\newcommand{\filepath}[1]{{\small\ttfamily\color{bpsmuted}\nolinkurl{#1}}}
\newcommand{\refbadge}{{\small\colorbox{bpsenh}{\textcolor{white}{\bfseries\ PROPOSED REFINEMENT \ }}\quad\itshape\color{bpsenh}Awaiting expert evaluation. Not yet applied to the pipeline.\par\smallskip}}
\newcommand{\schemeheader}[3]{%
  \begin{center}
    {\LARGE \bfseries\color{bpsink} #1}\\[0.45em]
    {\large #2}\\[0.35em]
    {\normalsize \itshape\color{bpsmuted} #3}
  \end{center}
  \vspace{0.4em}\hrule\vspace{0.9em}
}
\newcommand{\statusnote}[2]{%
  \begin{center}
  \fcolorbox{bpsline}{bpswarnbg}{%
    \parbox{0.94\linewidth}{\small\color{bpswarnink}\textbf{#1.}\ #2}%
  }
  \end{center}
  \vspace{0.6em}
}


\begin{document}
\schemeheader
  {Scheme 4}
  {Stage 3 Retrieval and Manual Relevance Triage Scheme}
  {Full-text availability, retrieval need, and adjudication checklist}

\statusnote{DRAFT FOR EXPERT EVALUATION}{These coding schemes are a working draft circulated for expert evaluation. They have not been applied to a final review corpus. The current manuscript is a test run that exercised an earlier, coarser generation of these schemes. The workflow itself has since been validated end to end in two cross-provider test runs, in which three large language models from three different providers applied the abstract-level and the full-text scheme independently and their agreement was quantified per coded field. The full run on the review corpus is deliberately held until this evaluation is complete.}

\section*{At a Glance}
\begin{description}
\item[Workflow position] Bridge between Stage 2 coding and Stage 3 full-text coding.
\item[Operational mode] Automated candidate manifest plus human checklist completion.
\item[Unit of analysis] One Stage 3 candidate record and its retrieval state.
\item[Provenance basis] Manifest-generation logic in stage3\_prep.py and the generated manual relevance checklist.
\item[Research questions] Retrieval and adjudication gate that protects RQ2 and RQ3 corpus quality
\item[Tagline] Automated candidate manifest plus human checklist completion
\end{description}

\section*{Canonical Source Paths}
\begin{itemize}
\item \filepath{src/03_pipeline/bps_review/extraction/stage3_prep.py}
\item \filepath{src/09_review_stages/04_extraction/forms/stage3_manual_relevance_checklist.csv}
\item \filepath{src/09_review_stages/04_extraction/outputs/stage3_candidate_manifest.csv}
\item \filepath{src/09_review_stages/04_extraction/outputs/stage3_manual_fulltext_queue.csv}
\item \filepath{src/09_review_stages/04_extraction/outputs/stage3_retrieval_validation.csv}
\end{itemize}

\section*{Purpose}
This scheme governs the transition from Stage 2 abstract coding to Stage 3 full-text work. It standardizes which candidate reviews need manual retrieval, which records require manual relevance adjudication, and how retrieval status, risk signals, and reviewer decisions are recorded before deep coding begins.

It is not a logistics file. It is a standardized adjudication framework that decides whether a review can enter Stage 3, whether the retrieved text is adequate, and how problematic or ambiguous records are escalated for human judgment.


\section*{Manifest Fields}
{\small\itshape Automated fields written per candidate.}\par\smallskip
\begin{description}
\item[\texttt{fulltext\_status}] Operational retrieval status of the candidate.
  \begin{itemize}[leftmargin=1.3em,itemsep=2pt,topsep=2pt]
  \item \textbf{\texttt{manual\_retrieval\_required}} No machine-available full text was secured; human retrieval is needed.
  \item \textbf{\texttt{pmc\_open\_available\_not\_cached}} A PMCID exists but the file was not yet cached.
  \item \textbf{\texttt{pmc\_fulltext\_cached}} A PMC full text is already cached locally.
  \item \textbf{\texttt{pmc\_fulltext\_fetched}} A PMC full text was retrieved during the current run.
  \item \textbf{\texttt{pmc\_linked\_fetch\_failed}} A PMCID link existed but retrieval or XML parsing failed.
  \item \textbf{\texttt{pmc\_fulltext\_low\_content\_manual\_check}} Text was fetched but appears too short and needs manual inspection.
  \end{itemize}
\item[\texttt{review\_track}] Which review pool the candidate belongs to, carried from Scheme 2. Retrieval and triage logic itself is uniform across both tracks.
  \begin{itemize}[leftmargin=1.3em,itemsep=2pt,topsep=2pt]
  \item \textbf{\texttt{musculoskeletal}} Musculoskeletal review pool.
  \item \textbf{\texttt{neuropathic}} Neuropathic review pool.
  \item \textbf{\texttt{both}} In both pools.
  \end{itemize}
\item[\texttt{retrieval\_source}] Source of the PMCID or full-text link (existing metadata, PubMed elink, or Europe PMC).
\item[\texttt{fulltext\_word\_count}] Word count of cached text, used to detect low-content full texts. \textit{(free text)}
\item[\texttt{manual\_retrieval\_needed}] Whether a human must retrieve the text.
  \begin{itemize}[leftmargin=1.3em,itemsep=2pt,topsep=2pt]
  \item \textbf{\texttt{yes}} Retrieval required.
  \item \textbf{\texttt{no}} Machine text is available.
  \end{itemize}
\item[\texttt{manual\_relevance\_priority}] Adjudication urgency.
  \begin{itemize}[leftmargin=1.3em,itemsep=2pt,topsep=2pt]
  \item \textbf{\texttt{high}} Withdrawal or missing pain focus.
  \item \textbf{\texttt{medium}} Other concern flags, or retrieval required with otherwise low signal.
  \item \textbf{\texttt{low}} No concern flags.
  \end{itemize}
\item[\texttt{manual\_relevance\_flags}] Pipe-delimited signal list describing why adjudication is needed. See the signal logic below. \textit{(free text)}
\item[\texttt{osf\_manual\_adjudication\_required}] Currently set to yes for the checklist workflow.
\item[\texttt{cached\_xml\_path / cached\_text\_path}] Relative cache paths for machine-fetched PMC files.
\item[\texttt{pubmed\_url}] Direct URL to the source record when a PMID is available.
\end{description}

\section*{Manual Relevance Signal Logic}
{\small\itshape Each signal is derived from title and abstract and drives priority.}\par\smallskip
\begin{description}
\item[\texttt{withdrawn\_or\_retracted\_signal}] Assigned when the title or abstract indicates a withdrawn or retracted record. Forces high priority.
\item[\texttt{pain\_focus\_not\_explicit}] Assigned when pain is not explicit in title or abstract. Forces high priority.
\item[\texttt{chronicity\_not\_explicit}] Assigned when chronic, persistent, or long-term framing is absent.
\item[\texttt{review\_design\_unclear}] Assigned when review design is missing or too vague.
\end{description}

\section*{Reviewer Checklist Fields}
{\small\itshape Human adjudication fields completed against the manifest.}\par\smallskip
\begin{description}
\item[\texttt{reviewer\_decision}] Human reviewer decision after checking the candidate. \textit{(free text)}
\item[\texttt{reviewer\_notes}] Notes supporting the reviewer decision. \textit{(free text)}
\item[\texttt{adjudication\_decision}] Final adjudicated status after disagreement resolution. \textit{(free text)}
\item[\texttt{adjudication\_notes}] Adjudication rationale. \textit{(free text)}
\end{description}

\section*{Primary Outputs Using This Scheme}
\begin{itemize}
\item \filepath{src/09_review_stages/04_extraction/outputs/stage3_candidate_manifest.csv}
\item \filepath{src/09_review_stages/04_extraction/forms/stage3_manual_relevance_checklist.csv}
\item \filepath{src/09_review_stages/04_extraction/outputs/stage3_manual_fulltext_queue.csv}
\item \filepath{src/09_review_stages/04_extraction/outputs/stage3_retrieval_validation.csv}
\end{itemize}

\end{document}
